
%%%%%%%%%%%%%%%%%%%%%%%%%%%%%%%%%%%%%%%%%%%%%%%%%%%%%%%%%%%%%%%%%%%%%%%%%%%%%%%%%%
\subsection{ Gauge  transformations}
\label{sec:velocity_gauge}

In this section we discuss traansformation og Time-Dependent Kohn-Sham  equation  under differet gauge representation within Coulomb gauge condtion.
The srtarting point are  general Eqs. [ A'->A+ '  ; Phi'->Phi +]   and  the requirement that physically measurable quantities such  mechanical momentum, velocities
 and  physicall measurable  operator observable remain unchanged 
The  invarance of E and B is directly guaranted by   [XXX] 


For finite systems, the relevant observable is the
time-dependent dipole moment $\vb d(t) = \int \vb r\,
\rho(\vb r,t)\, d\vb r$, from which the polarizability
and absorption cross section follow by Fourier transform.
For periodic systems, however, the scalar-potential form
$\Phi = -\vb E(t)\cdot\vb r$ is not compatible with
periodic boundary conditions, because the position
operator $\vb r$ is discontinuous at the cell boundary.
This motivates the use of the velocity-gauge formulation,
in which the external field enters through a spatially
uniform vector potential $\vb A(t)$ and the scalar
potential is set to zero.

The time-dependent Kohn--Sham equation then takes the form
\begin{equation}
  \imath \pdv{}{t} \Psi_j(\vb r,t) = \left[
    \frac{1}{2m}\left(\vb{\hat p}
    - \frac{q_e}{c}\vb A(t)\right)^2
    + \hat V_{\mathrm{KS}}[\rho](t)
  \right] \Psi_j(\vb r,t),
  \label{eq:tdks_velocity_gauge}
\end{equation}
where the two gauges are related by the Göppert-Mayer
transformation, discussed in detail in
Sec.~\ref{sec:velocity_gauge}.
In the velocity gauge, the position operator no longer
appears in the Hamiltonian and all terms are compatible
with periodic boundary conditions.
Correspondingly, the physical observable used to extract
the optical response is the time-dependent current density
$\vb j(\vb r,t)$ rather than the dipole moment, since
the current is well defined under periodic boundary
conditions whereas $\vb d(t)$ is not.







